%% kapa: exam=1 round=27 year=2016 period=1 form=A subject=economics subject_ko=경제학원론 seq=35 number=75
%% source: Q-Net 감정평가사 기출문제 (cst00309) artlSeq=5169402 file=extracted/1008/5169402/제27회 감정평가사 1차 문제/제27회 감정평가사 1교시문제(A형).pdf
%% conversion: kordoc 4.24.51 -> tools/gen1.py (auto); review=proofread-sonnet; flags=none
\begin{question}{75}
사과와 오렌지만 생산하는 A국의 생산량과 가격이 다음과 같을 때 2014년 대비 2015년의 $GDP$ 디플레이터로 계산한 물가상승률은 얼마인가? (단, 2014년을 기준년도로 한다.)
\par{\small\centering
\begin{tabular}{|c|c|c|c|c|}\hline
\multirow{2}{*}{연도} & \multicolumn{2}{c|}{사과} & \multicolumn{2}{c|}{오렌지} \\ \cline{2-5}
 & \makebox[1.8cm]{수량} & \makebox[1.8cm]{가격} & \makebox[1.8cm]{수량} & \makebox[1.8cm]{가격} \\ \hline
2014 & 5 & 2 & 30 & 1 \\ \hline
2015 & 10 & 3 & 20 & 1 \\ \hline
\end{tabular}\par}
\begin{choices}
  \item 20\%
  \item 25\%
  \item 35\%
  \item 45\%
  \item 50\%
\end{choices}
\end{question}
\answer{2}
